% Round 22 / geometry branch. Candidate; independent review required.
\section{Two surface factors, actual triple strata, and the full fourfold gate}
\label{r22ds:section}

\paragraph{Scope.}
This section gives a new local presentation beyond the one-curve-factor
construction MX21.  Its generic fourfold is an ample divisor in
$S\times T\times\mathbf P^1$.  The central fiber has two curve
factors and several point factors; its actual triple strata are products
of the two curves.  The proof uses the full numerical middle space of
each component.  For the point components $S\times T$, that positivity
is an explicit hypothesis, with an unconditional elliptic-square
instance proved below from the frozen FP21 criterion.  Positivity of
an arbitrary product of two surfaces is not asserted.

All numerical groups and intersection degrees below are native to the
indicated finite field.  The symbol $N^a$ means rational cycles modulo
numerical equivalence.  Numerical duality is the perfect pairing between
$N^a$ and $N^{\dim-a}$, not a comparison with an unproved cohomological
quotient.  The final Green output has the frozen ordinary actual
$\delta$-form meaning.

\subsection{The supplied local datum}

Let $R=\mathcal O_K$ be the valuation ring of the fixed finite valued
extension $K/\mathbf Q_p$, with native uniformizer $\pi$ and finite
residue field $k$.  Let $\mathcal S,\mathcal T$ be smooth projective
geometrically integral surface schemes over $R$.  Write $S,T$ for
their special fibers.  Choose relatively ample line bundles
$\mathcal A,\mathcal B$ and sections $f,g$ whose special zero divisors
$C\subset S,D\subset T$ are smooth geometrically integral curves.
Choose $d\ge2$ linear forms $l_1,\ldots,l_d$ on $\mathbf P^1_R$
with pairwise distinct native special zeroes $p_1,\ldots,p_d$.
Put
\[
 \mathcal Y=\mathcal S\times_R\mathcal T\times_R\mathbf P^1_R,
 \qquad \mathcal L=\mathcal A\boxtimes\mathcal B\boxtimes\mathcal O(d),
 \qquad P=fg\prod_{j=1}^d l_j.
 \tag{DS22.1}\label{r22ds:datum}
\]
The line bundles of the individual factors need not agree.  On the
special fiber, order the factor divisors as
\[
 H_0=C\times T\times\mathbf P^1,\quad
 H_1=S\times D\times\mathbf P^1,\quad
 H_{j+1}=S\times T\times\{p_j\}\quad(1\le j\le d).
 \tag{DS22.2}\label{r22ds:factors}
\]
For a nonempty index set $I$, let $H_I=\bigcap_{i\in I}H_i$.
No two point factors meet, so $H_I$ is empty for $|I|\ge4$.
Every nonempty $H_I$ is smooth and geometrically integral.

Choose $Q\in H^0(\mathcal Y,\mathcal L)$ with these precise properties:
its special section cuts every nonempty $H_I$ transversely; every
positive-dimensional cut is geometrically integral; every
zero-dimensional cut is finite etale.  Empty cuts are allowed.
Assume also that
\[
                    X=V(P+\pi Q)\subset\mathcal Y_K
 \tag{DS22.3}\label{r22ds:X}
\]
is smooth and geometrically integral.  These are supplied hypotheses,
not a claim that arbitrary coefficients satisfy them.  For the ample
elliptic product data used in DG22, simultaneous Bertini gives such
residue sections after one finite unramified preparation.  The global
argument there constructs the smooth generic members separately.

\subsection{Actual resolution and its complete strata}

\begin{proposition}[Ordered resolution, DS22]
\label{r22ds:resolution}
Blow up the ambient $\mathcal Y$ successively in the strict transforms
of $H_0,H_1,\ldots,H_d$ as closed subschemes of its special fiber, and
take the strict transform of the flat closure of $X$.
The result $\mathcal X$ is projective, regular and reduced strictly
semistable over $R$.  Its components are the $d+2$ continuing factor
components $D_0,\ldots,D_{d+1}$ and no others.  Every nonempty
stratum $D_I$ is the iterated smooth blowup of $H_I$ in the centers
\[
              Z_{I,j}=H_I\cap H_j\cap V(\bar Q),\qquad j>\max I,
 \tag{DS22.4}\label{r22ds:centers}
\]
in decreasing order of $j$.  Every nonempty center has codimension
two in $H_I$, and its transform at the moment it is blown up is
isomorphic to that original center.  In particular the full double
and triple list is
\[
\begin{aligned}
 D_{01}&=\operatorname{Bl}_{\coprod_{j=1}^d Z_j}
                          (C\times D\times\mathbf P^1),
 &Z_j&=(C\times D\times\{p_j\})\cap V(\bar Q),\\
 D_{0,j+1}&=C\times T,
 &D_{1,j+1}&=S\times D,\\
 D_{0,1,j+1}&=C\times D.
\end{aligned}
 \tag{DS22.5}\label{r22ds:explicitstrata}
\]
Here the $Z_j$ are disjoint smooth curves when nonempty.  There are
genuine triple strata and no quadruple strata.  Each component is a
sequence of smooth surface blowups of either
$C\times T\times\mathbf P^1$, $S\times D\times\mathbf P^1$,
or $S\times T$; the point components $D_{j+1}=S\times T$ are
unchanged.
\end{proposition}

\begin{proof}
The actual local calculation is the same ordered-factor calculation
as MX21, but we give it to fix the centers and eliminate any appeal
to an unspecified resolution.  At a point on $s$ central factors,
their defining functions are distinct members of a regular relative
parameter system.  Here $s\le3$.  If $\bar Q$ is a unit, the equation
is $x_1\cdots x_s=\pi u$ with $u$ a unit; it is already strictly
semistable and the factor ideals become principal.  If $\bar Q=0$,
the transverse cut supplies one additional parameter, and the local
equation, after a harmless unit change, is
\[
                          x_1\cdots x_s=\pi y.
 \tag{DS22.6}\label{r22ds:local}
\]
The two saturated charts of the first blowup in $(\pi,x_1)$ are
$\pi=x_1t_1$, $x_2\cdots x_s=t_1y$, and
$x_1=\pi u_1$, $y=u_1x_2\cdots x_s$.  Continuing in the first
chart gives the terminal charts
\[
 \pi=x_1\cdots x_{a-1}t_{a-1},\qquad
 x_a=t_{a-1}u_a,\qquad y=u_ax_{a+1}\cdots x_s.
 \tag{DS22.7}\label{r22ds:charts}
\]
After eliminating the two displayed variables the remaining local
ring is regular and $\pi$ is a square-free product.  The final
continuing chart has the same property.  The vertical factors in
these charts are precisely the original factor components.  Thus
there are no additional components or hidden multiplicities.
Projectivity follows from the actual ambient blowups.  The closure
and its transform are integral, because the assumed geometrically
integral generic fiber is dense; the displayed charts and the
relative Cartier equation also verify flatness.

On an earlier factor the modification induced by a later factor is
the two-generator ideal $(\bar f_j,\bar Q)$; on the later factor
itself that step is Cartier and has no effect.  Restricting the two
Rees charts to an intersection indexed by $I$ leaves precisely the
ideals with $j>\max I$, in the reverse order displayed in
\eqref{r22ds:centers}.  Local unit changes preserve these ideals,
so the chart calculations glue.  Transversality makes each center
smooth of codimension two.  A preceding center restricts to a
subsequent center by the ideal of the remaining factor, since
$\bar Q$ already vanishes there.  It is Cartier, so that blowup
does not change the subsequent center.  The possible intersection
of two different point factors is empty.  This proves the center
assertions without a hidden normal-crossings simplification.

For $I=\{0,1\}$ the only remaining factors are the $d$ distinct
points, giving the disjoint curves $Z_j$.  For $I=\{0,j+1\}$ or
$\{1,j+1\}$ there is no larger intersecting point factor and no
remaining earlier surface factor, so these doubles are unchanged.
The triples are likewise unchanged.  This proves the entire list
\eqref{r22ds:explicitstrata}.  On a single factor all centers are
surfaces, and on a point component there are no centers.  These
observations also prove the component assertion.
\end{proof}

\subsection{A full hyperbolic middle form independent of standard D}

\begin{lemma}[Threefold times a projective line]
\label{r22ds:hyperbolic}
Let $W/k$ be any smooth projective geometrically integral threefold.
On the full native numerical middle group of $W\times\mathbf P^1$
the intersection form has signature
$(\rho(W),\rho(W))$.  Its negative index is
$\rho(W\times\mathbf P^1)-1$.  For every ample class $H$ on this
fourfold the primitive middle group
$\ker(H:N^2\to N^3)$ is positive definite.
\end{lemma}

\begin{proof}
The actual projective-bundle formula and numerical duality give
\[
 N^2(W\times\mathbf P^1)
            =N^2(W)\oplus N^1(W)h,
 \qquad h^2=0.
 \tag{DS22.8}\label{r22ds:hyperbolicform}
\]
Each summand is isotropic.  Their cross pairing is precisely the
perfect native numerical pairing $N^2(W)\times N^1(W)\to\mathbf Q$.
Both summands therefore have dimension $\rho(W)$, and the form is
hyperbolic.  Moreover $\rho(W\times\mathbf P^1)=\rho(W)+1$.
This calculation uses all numerical cycles, including classes not
known to arise from an algebraic basis of geometric cohomology.

Divisor Hodge index for the ample class $H$ makes the pairing on
$HN^1$ nondegenerate of signature $(1,\rho(W))$; in particular
multiplication by $H$ on $N^1$ is injective.  Its orthogonal
complement in $N^2$ is $\ker(H:N^2\to N^3)$ by numerical duality.
All negative directions already lie in $HN^1$.  The full form is
nondegenerate, so that orthogonal complement is positive definite.
No assertion about the Lefschetz standard conjecture for an
arbitrary threefold was used.
\end{proof}

Use the following explicit condition on the two surface fibers:
\[
 b^-\bigl(N^2(S\times T)_{\mathbf R},\deg(xy)\bigr)
                         =\rho(S\times T)-1.
 \tag{DS22.9}\label{r22ds:MP}
\]
When an all-finite-extension assertion is wanted, require this same
condition after every finite extension of $k$, for the whole new
native numerical group there.

\begin{proposition}[Full component positivity]
\label{r22ds:positivity}
Assume \eqref{r22ds:MP}.  Every component $D_i$ of the actual model
has negative middle index $\rho(D_i)-1$.  For the restrictions of
any one genuinely relatively ample line bundle on $\mathcal X$,
all component primitive middle forms are positive definite.
\end{proposition}

\begin{proof}
Lemma~\ref{r22ds:hyperbolic} applies to the original first two
components, taking $W=C\times T$ or $W=S\times D$.
Condition \eqref{r22ds:MP} applies to each unchanged point component.
For a smooth surface center $Z$ in a smooth fourfold $M$, write
$q:\widetilde M=\operatorname{Bl}_ZM\to M$, $j:F\hookrightarrow
\widetilde M$, $p:F\to Z$.  The actual Chow blowup decomposition
and its native degree pairing give
\[
 N^2(\widetilde M)
           =q^*N^2(M)\ \perp\ j_*p^*N^1(Z),\qquad
 (j_*p^*a,j_*p^*b)=-\deg_Z(ab).
 \tag{DS22.10}\label{r22ds:blowupsignature}
\]
The mixed term is a codimension-three intersection on the surface
$Z$ and is zero.  The exceptional formula follows from
$j^*j_*=-\xi$ and $p_*\xi=1$.  Both summand pairings are
nondegenerate, which also justifies passing from the Chow formula
to the numerical direct sum.  Surface divisor Hodge index shows
that each geometrically connected surface center adds exactly one
negative direction, while the Picard rank adds one.  The supplied
integrality hypotheses ensure the centers here are geometrically
integral; disjoint connected centers could equivalently be handled
componentwise with their native degrees.  Induction through the
actual centers of \eqref{r22ds:centers} proves the index formula.

Finally, for any ample restriction $H_i$, divisor Hodge index on
$H_iN^1(D_i)$ already accounts for the entire negative index.
Its orthogonal complement is the primitive middle group and is
positive definite, exactly as in the preceding lemma.  A merely
nef pullback class has not been substituted for the required ample
class.
\end{proof}

\begin{lemma}[The complete native curve gate on the doubles]
\label{r22ds:CD}
Every double in \eqref{r22ds:explicitstrata} has the CD17 property:
an actual native numerically zero rational curve has zero geometric
$\ell$-adic class, for every $\ell\ne\operatorname{char}k$.
The assertion holds anew after every finite constant extension.
\end{lemma}

\begin{proof}
The original doubles $C\times T$, $S\times D$, and
$C\times D\times\mathbf P^1$ are products of smooth curves and
surfaces.  The frozen SP19 product comparison applies to their
entire curve groups.  Its precise classical input is algebraicity
of inverse Lefschetz in degree two for products of curves and
surfaces, not algebraicity of all $H^2$.  The only transformed
double is the blowup $D_{01}$ in the smooth curves $Z_j$.
Frozen CD17 smooth-blowup comparison, or SP19 birational ascent,
therefore gives the assertion for it.  The same proofs apply to
the whole new native curve group over each finite constant field;
no new cycle is assumed to descend to $k$.
\end{proof}

\begin{theorem}[Full $U_2$ and complete local G, DS22]
\label{r22ds:green}
For every datum \eqref{r22ds:datum}--\eqref{r22ds:X} satisfying
\eqref{r22ds:MP}, the actual reduced strictly semistable model has
\[
 U_2=\frac{\ker\rho_2\cap\operatorname{im}\sigma_2}
                 {\operatorname{im}(\sigma_2\rho_1)}=0.
 \tag{DS22.11}\label{r22ds:U}
\]
Its generic fourfold has full strong G in every codimension.
If \eqref{r22ds:MP} holds after every finite constant extension,
then for every finite valued extension $L/K$ the same complete
assertion holds for every new specified geometrically homologically
trivial rational cycle on $X_L$.
\end{theorem}

\begin{proof}
Choose an actually relatively ample line bundle on $\mathcal X$.
Frozen PU18 identifies the full $U_2$ with its unique primitive
representative subspace in $\ker\rho_2\cap\operatorname{im}\sigma_2$
and proves that subspace totally isotropic.  The direct sum of the
component primitive middle forms is positive definite by
Proposition~\ref{r22ds:positivity}; hence that subspace is zero.
The model has genuine triples throughout this argument.  Neither
$\ker A=\ker\rho$ nor $K_{\rm comp}=\ker A$ is claimed.

Lemma~\ref{r22ds:CD} closes the actual representative-lifting gate
of CL17.  The frozen exact detector, equivalently the MZ20
consequence of $U_2=0$, kills the remaining cycle-visible obstruction
for every original geometrically homologically trivial codimension
two Chow class.  The frozen divisor, one-cycle, and zero-cycle
theorems supply all other fourfold codimensions.  Thus VT17 gives,
for every specified $Z$, an actual ordinary-topology form
\[
 \gamma_Z\in B^{r-1,r-1}
                  (X^{\rm an}\setminus|Z|^{\rm an})
 \quad\text{with}\quad
 \int_{X^{\rm an}\setminus|Z|^{\rm an}}^*
       \gamma_Z\wedge d'd''\eta=-\int_{Z^{\rm an}}^*\eta
 \tag{DS22.12}\label{r22ds:actualG}
\]
for every actual compactly supported complementary test $\eta$.
This is the exact original support and the full test space.

For each finite unramified extension, the same smooth strata and
the ordered charts base change, and the assumed positivity is for
the whole new native numerical group.  Reapply the argument to
every new cycle over that field.  The model is reduced strictly
semistable, so frozen RI18, Corollary \texttt{r18ri:all-finite},
propagates the result from all unramified fields to every finite
valued extension.  Its norm argument and finite actual trace retain
the full degrees, including ramification and residue degrees, and
the extension of the fixed original absolute value.  The last
application of VT17 is made to the final specified cycle.  No
arbitrary proper actual pushforward or finite-on-complement
alteration assertion enters.
\end{proof}

\subsection{A concrete unconditional positivity input}

\begin{corollary}[Ordinary elliptic squares at five]
\label{r22ds:elliptic}
Let $E/\mathbf F_5$ be $y^2=x^3-x$, and let $S=T=E^2$ after any
finite constant extension.  Condition \eqref{r22ds:MP} holds after
every further finite constant extension.  Consequently every local
smoothing datum above with these good special surfaces has the
complete all-finite-extension Green property of
Theorem~\ref{r22ds:green}.
\end{corollary}

\begin{proof}
Counting the five possible $x$ coordinates gives $\#E(\mathbf F_5)=8$.
The Frobenius eigenvalues on $H^1(E)$ are
$\alpha=-1+2i$, $\beta=-1-2i$, with $\alpha\beta=5$.
The two factor divisors, the diagonal, and the graph of the
automorphism $(x,y)\mapsto(-x,2y)$ of $E^2$ have an intersection
matrix of determinant $-4$.  Thus they are four independent native
divisor classes.  On $H^2(E^2)$ the remaining two eigenvalues are
$\alpha^2=-3-4i$ and $\beta^2=-3+4i$.

The ratio $\alpha/\beta=(-3-4i)/5$ is not a root of unity: the only
roots of unity in $\mathbf Q(i)$ are $1,-1,i,-i$.
Therefore $\alpha^2/5$, $\beta^2/5$, and
$\alpha^2/\beta^2=(\alpha/\beta)^2$ are not roots of unity.
A geometric divisor is defined over a finite constant field and
has a Tate eigenvalue after that extension.  Hence the two remaining
directions never become divisors, and the four displayed classes
are the full geometric divisor space over every constant field.
The actual full geometric divisor projector in FP21 is therefore
native already over $\mathbf F_5$.

Frobenius on the transcendental surface block has the two distinct
eigenvalues $\alpha^2,\beta^2$, and every positive power still has
distinct eigenvalues.  Frozen FP21, Corollary
\texttt{r21fp:simplespectrum}, now gives the full numerical middle
signature of $S\times S$ and its all-finite-constant-field assertion.
This is exactly \eqref{r22ds:MP}, for all actual new cycles.  It
does not discard a possible $t_2\otimes t_2$ block: FP21 identifies
that whole block by the actual Frobenius correspondence algebra.
\end{proof}

\paragraph{Dependencies and the remaining boundary.}
The geometric local calculation extends the explicitly checked
ordered-factor calculation of MX21.  The full positive-product
input in the displayed elliptic case is the frozen FP21 simple
all-powers Frobenius theorem.  The threefold-times-line hyperbolic
calculation restates the already closed rank-two projective-bundle
case of PC16, itself retaining its R7 provenance; it is not a new
positivity result.  The new content is its activation on the actual
two-surface mixed models, together with their complete strata and
cycle domain.  The final
interfaces are PU18, CD17/SP19, CL17, MZ20, RI18, VT17 and the
already closed other-codimension Green theorems.  No old analytical
gate is reopened.  For arbitrary good surfaces $S,T$, condition
\eqref{r22ds:MP} remains an explicit hypothesis; this section does
not prove general G or arbitrary-variety complete heights.
